% Second pass (Claude, 2026-11-10): first pass checked against the archive images; changes in LM_C3-568-586_changes.md.
\scanpage{568}
% source-position: top, lines 1--2
immediate from (77). \deleted{It therefore remains to}\inserted{We thus have}
\deleted{\unsure{d}}

% source-position: upper, display (79)
\[
(79)\qquad
\diagram{\begin{tikzcd}
S_0\Gamma_{\rho,\sigma} \arrow[r,hook] & S\unsure{\mathfrak{G}}_{\rho,\sigma}^{\mathrm{eff}}
\end{tikzcd}}
\]
% source-position: upper, lines 3--5
it remains to see that this map is surjective ---
we have done what was needed for. Let us first note
that the composite
\[
\diagram{\begin{tikzcd}
S_0\Gamma_{\rho,\sigma} \arrow[r,hook] & S\unsure{\mathfrak{G}}_{\rho,\sigma}^{\mathrm{eff}} \arrow[r,"\text{epimorphism}"] & \Sigma_{\rho,\sigma}^{\mathrm{eff}}\bigl( \arrow[r,"\sim"] & Z_{\rho,\sigma}\bigr)
\end{tikzcd}}
\]
\deleted{\(\diagram{\begin{tikzcd}(u,a,b)\arrow[r,mapsto]&u\end{tikzcd}}\)}
% source-position: upper-middle, lines 6--7
is none other, by (77) and (71) than
\deleted{\(\diagram{\begin{tikzcd}
(\unsure{\alpha},\unsure{\beta},\unsure{\mu})\arrow[r,mapsto]&
(\unsure{[\alpha,\rho]},\unsure{[\beta,\sigma]})
\end{tikzcd}}\)}
\[
(80)\qquad
\diagram{\begin{tikzcd}
x=(u,a,b)\arrow[r,mapsto]& (\xi,\eta,\mu,\varepsilon',\varepsilon)
\end{tikzcd}}
\quad
\left\{\begin{array}{ll}
\xi=[u,\rho]&\text{thus }u(\rho)=\xi\rho\\
\eta=[u,\sigma]&\text{thus }u(\sigma)=\eta\sigma\\
\mu=\chi_{S\Gamma}(x)&\\
\varepsilon'=\varepsilon'(x)&\\
\varepsilon=\varepsilon(x)&
\end{array}\right.
\]
% source-position: middle, lines 8--10
thus its \deleted{image} composite with
\(\diagram{\begin{tikzcd}\Sigma_{\rho,\sigma}^{\mathrm{eff}}\arrow[r,"\sim"]&Z_{\rho,\sigma}\end{tikzcd}}\)
is none other than the canonical projection
\[
\diagram{\begin{tikzcd}
S_0\Gamma_{\rho,\sigma}\arrow[r]&Z_{\rho,\sigma}\\
(u,a,b)\arrow[r,mapsto]&u
\end{tikzcd}}
\]
which is surjective, with kernel
\(SN_\rho^*\times SN_\sigma^*=SZ_\rho\times SZ_\sigma\).

% source-position: lower-middle, lines 11--12
Let us consider on the other hand \deleted{\unsure{the operation}}
\deleted{It remains to see that}
\[
(81)\quad
\left\{\begin{array}{l}
\text{operation of }SN_\rho\times SN_\sigma\\
\text{on }\deleted{S}\unsure{\mathfrak{G}}_{\rho,\sigma}
\subset\mathcal{G}\times\mathcal{G}\times\widehat{\mathbb Z}^{*}\\
\hfill\Vert\\
\{(\alpha,\beta,\mu)\mid[\alpha,\rho]=[\beta,\sigma]\varepsilon_1^{\mu-1}\}
\end{array}\right.
\quad
\begin{array}{l}
(a_0,b_0)\in S\deleted{\unsure{N}}Z_\rho^{\unsure{*}}\times S\deleted{\unsure{N}}Z_\sigma^{\unsure{*}}\\
\inserted{(=SZ_\rho\times SZ_\sigma)}\\
\text{operates by}\\
\diagram{\begin{tikzcd}(\alpha,\beta,\mu)\arrow[r,mapsto]&(\alpha a_0,\beta b_0,\mu)\end{tikzcd}}
\end{array},
\]
% source-position: lower, lines 13--17
It is clear that it operates \deleted{faithfully} freely, and
that the quotient of \(\unsure{\mathfrak{G}}_{\rho,\sigma}\) by this action
is identified with \(\Sigma_{\rho,\sigma}\), thus we have the result
idem for \(\unsure{\mathfrak{G}}_{\rho,\sigma}^{\unsure{*}}\) and
\(\Sigma_{\rho,\sigma}^{\unsure{*}}\) \deleted{\illegible{1}}. Now, for
% source-position: bottom, overwritten display (82), lines 18--20
\deleted{(82) for \(\varepsilon,\varepsilon'\)
\inserted{\(\varepsilon=(-1)^\nu,\quad\varepsilon'=(-1)^{\nu'}\)} fixed,
the action (81) of \(SZ'_\rho\times SZ'_\sigma\) corresponds
in \(S_0\Gamma_{\rho,\sigma}\) to the \deleted{translation}\inserted{operation}
\unsure{\((u,a,b)\)}
\(\diagram{\begin{tikzcd}
(u,a,b)\arrow[r,mapsto]&
\bigl(u,\unsure{\varepsilon_\rho^{\tau'}(a_0^{-1})}a,
\unsure{\varepsilon_\sigma^{\tau'}(b_0^{-1})}b\bigr)
\end{tikzcd}}\)}
\mspage{567}

% Printed archive footer outside the manuscript frame (not manuscript):
% 567
% https://grothendieck.umontpellier.fr
% Université de Montpellier, UMR CNRS 5149, Case courrier 051, Place Eugène Bataillon
% 34095 Montpellier cedex 5 - France

\scanpage{569}
% source-position: top edge, marginal calculation
\margin{\deleted{\unsure{i.e.}\ \(\unsure{\varepsilon_\rho^{\tau'}}\,
\unsure{(a)}=\tau'^{\nu'}\)}}
\mspage{588}
% source-position: top line 1, cancelled
\deleted{\unsure{it is thus} (\illegible{2} \unsure{the set}}
% source-position: upper display (82)
\[
(82)\qquad \varepsilon=(-1)^\nu,\quad\varepsilon'=(-1)^{\nu'}
\]
fixed, the preceding operation of \((a_0,b_0)\) corresponds,
under the map (76), to the operation
\[
(83)\quad
\diagram{\begin{tikzcd}
(u,a,b)\arrow[r]&\bigl(u,\tau^{\nu'}(a_0^{-1})a,\tau^\nu(b_0^{-1})b\bigr)
\end{tikzcd}},
\]
% source-position: upper-middle, lines 5--8
i.e. to left translation by
\((\tau^{\nu'},\tau^\nu)\underbrace{(a_0^{-1},b_0^{-1})}_{\inserted{(a_0,b_0)^{-1}}}\).
Whence surjectivity, and the desired isomorphism
\[
(84)\qquad
\diagram{\begin{tikzcd}
S_0\Gamma_{\rho,\sigma}\arrow[r,"\sim"]&S\unsure{\mathfrak{G}}_{\rho,\sigma}^{\mathrm{eff}}
\end{tikzcd}}
\quad\text{(cf. 76)}
\]
% source-position: middle, lines 9--12
The left translation by \(\ell\in L_0\)
\inserted{in \(S_0\Gamma_{\rho,\sigma}\)} corresponds,
under the preceding bijection, to the \deleted{translation}\inserted{operation}
\[
(85)\qquad
\diagram{\begin{tikzcd}
(\alpha,\beta,\mu)\arrow[r,mapsto]&(\ell\alpha,\ell\beta,\mu)
\end{tikzcd}}
\]
which it was clear a priori stabilizes
\(S\unsure{\mathfrak{G}}_{\rho,\sigma}^{\mathrm{eff}}\).
Passing to the quotient, we thus find
\[
(86)\qquad
\diagram{\begin{tikzcd}
\Gamma_{\rho,\sigma}\arrow[r,"\sim"]&L_0\backslash S\unsure{\mathfrak{G}}_{\rho,\sigma}^{\mathrm{eff}}
\end{tikzcd}}.
\]

% source-position: lower, lines 13--14 and display (87)
Finally, let us consider
\[
(87)\quad
\left\{\begin{array}{l}
\diagram{\begin{tikzcd}[column sep=large]
S_0M_{\rho,\sigma}\arrow[r]&S\unsure{\mathfrak{G}}_{\rho,\sigma}^{\mathrm{eff}}\times\mathcal{G}^+
=\underbrace{\{\overset{\inserted{(\in\mathcal{G}\times\mathcal{G}\times\widehat{\mathbb Z}^{*}\times\mathcal{G}^{+}}}{(\alpha,\beta,\mu,f)}\mid}
_{\unsure{(}[\alpha,\rho]=[\beta,\sigma]\varepsilon_1^{\mu-1}\unsure{)}}\\
x=(u,a,b,U)\arrow[r,mapsto]&
\deleted{((\unsure{ua^{-1}\tau'},\unsure{ub^{-1}\tau'}),\unsure{uU^{-1}})}
\inserted{((\alpha,\beta,\mu),f)}
\end{tikzcd}}
\\
\hfill
\left\{\begin{array}{l}
\alpha=ua^{-1}\tau'^{\deleted{\unsure{\nu'}}\inserted{\unsure{\nu'_\rho(a)}}}\\
\beta=ub^{-1}\tau'^{\unsure{\nu_\sigma(b)}}\\
\mu=\chi_{SM}(x)\quad(=\chi_Z(u)=\chi_\rho(a)=\chi_\sigma(b)=\chi_G(U))\\
f=uU^{-1}
\end{array}\right.
\end{array}\right.
\]
% source-position: bottom line 15
This map is bijective (trivial by (84)).
\mspage{568}

% Printed archive footer outside the manuscript frame (not manuscript):
% 568
% https://grothendieck.umontpellier.fr
% Université de Montpellier, UMR CNRS 5149, Case courrier 051, Place Eugène Bataillon
% 34095 Montpellier cedex 5 - France

\scanpage{570}
\mspage{588}
% source-position: top heading
\textcircled{\(7^\circ\)}\quad\uline{Functorialities}

Let us consider
\[
(88)\qquad
\diagram{\begin{tikzcd}\mathcal{G}\arrow[r,"\varphi"]&\mathcal{G}'\end{tikzcd}}
\qquad\begin{array}{l}\text{surjective homomorphism of}\\\text{profinite groups}\end{array}
\]
% source-position: upper lines 3--5
let \deleted{\(\unsure{\rho,\sigma,\tau'}\)}
\(\rho_{\mathcal{G}'},\sigma_{\mathcal{G}'},\tau'_{\mathcal{G}'},\ldots\)
be the images of \(\rho=\rho_{\mathcal{G}},\sigma=\sigma_{\mathcal{G}},
\tau'=\tau'_{\mathcal{G}},\ldots\) in \(\mathcal{G}'\).

% source-position: upper-middle line 6
\uline{Proposition}\quad Equivalent conditions on \(\varphi\) :

a) The map induced by \(\varphi\)
\[
(89)\qquad
\diagram{\begin{tikzcd}
\Sigma_{\rho,\sigma}^{\mathcal{G}}\arrow[r]&\Sigma_{\rho,\sigma}^{\mathcal{G}'}
\end{tikzcd}}
\]
\deleted{\unsure{transforms}}\inserted{maps}
\(\Sigma_{\rho,\sigma}^{\mathcal{G}\,\mathrm{eff}}\)
into \(\Sigma_{\rho,\sigma}^{\mathcal{G}'\,\mathrm{eff}}\).

% source-position: middle lines 8--11
b) For every \(u\in\operatorname{Aut}\mathcal{G}^+\),
\unsure{which|that} normalizes \(L_0\),
and \unsure{such|such} that \(u(\deleted{\rho}L_\rho)\underset{\mathcal{G}^+}{\sim}L_\rho\),
\(u(L_\sigma)\underset{\mathcal{G}^+}{\sim}L_\sigma\),
there exists \(u'\in\operatorname{Aut}(\mathcal{G}'^+)\) which makes commutative
\[
(90)\qquad
\diagram{\begin{tikzcd}
\mathcal{G}^+\arrow[r,"\varphi"]\arrow[d,"u"']&\mathcal{G}'^+\arrow[d,"u'","{\unsure{-}}"']\\
\mathcal{G}^+\arrow[r,"\varphi"]&\mathcal{G}'^+
\end{tikzcd}},
\]
in other words, \(u\) \deleted{\unsure{maps}} \deleted{\unsure{stabilizes}}
\inserted{\unsure{respects}} \(\operatorname{Ker}\varphi\).

% source-position: middle-lower lines 12--15
b') idem, assuming only, in place of
\(u\) normalizing \(L_0\), that
\(u(L_0)\underset{\mathcal{G}^+}{\sim}L_0\)
(plus the other conditions)

\unsure{c|b}) \(\exists\) group homomorphism
\[
(91)\qquad
\varphi_Z:\quad
\diagram{\begin{tikzcd}
Z_{\rho,\sigma}^{\mathcal{G}}\arrow[r,"{\deleted{\unsure{\varphi_Z}}}"]&Z_{\rho,\sigma}^{\mathcal{G}'}
\end{tikzcd}}
\]
\deleted{\unsure{with} \(\unsure{u'}\)}
% source-position: lower lines 16--19
compatible with \(\chi,\varepsilon_Z,\varepsilon'_Z\) and such that
\(\forall u\in Z_{\rho,\sigma}^{\mathcal{G}}\), if
\(u'\) is its image we have commutativity in (90).
(\uline{N.B.} \(\varphi_Z\) \unsure{will be} \deleted{\illegible{1}}
unique \unsure{by|for} \unsure{\(\zeta\)|?}, and is compatible with \(i_{L_0,Z}\))
\margin{thus it induces
\(\diagram{\begin{tikzcd}
\gamma_{\rho,\sigma}^{\mathcal{G}}\arrow[r]&\gamma_{\rho,\sigma}^{\mathcal{G}'}
\end{tikzcd}}\)}

% source-position: bottom lines 20--23
c') \(\exists\) group homomorphism
\[
(92)\qquad
\varphi_{\deleted{Z}\inserted{G}}:\quad
\diagram{\begin{tikzcd}
G_{\rho,\sigma}^{\mathcal{G}}\arrow[r]&G_{\rho,\sigma}^{\mathcal{G}'}
\end{tikzcd}}
\]
compatible with \(\chi_G,\varepsilon_G,\varepsilon'_G\) and satisfying the
compatibility condition (90). (N.B. this \(\varphi_G\)
\mspage{569}

% Printed archive footer outside the manuscript frame (not manuscript):
% 569
% https://grothendieck.umontpellier.fr
% Université de Montpellier, UMR CNRS 5149, Case courrier 051, Place Eugène Bataillon
% 34095 Montpellier cedex 5 - France